\section{噪声影响维度的探索性因子分析降维}

噪声影响的14个题项之间存在较强相关性，若直接纳入后续模型会带来解释上的冗余。为提取结构清晰的潜在维度，本研究对上述题项进行探索性因子分析，将分散的题项信息压缩为少数几个相互独立的公因子，并为后续聚类分析与回归建模提供输入变量。

\subsection{适用性检验}

因子分析的前提是题项之间存在足够的\textbf{共同变异}。第四章的效度检验已针对这14个噪声影响题项完成KMO取样适切性检验与巴特利特球形检验，得到KMO取样适切性量数为0.815，巴特利特球形检验近似卡方值为2119.289（$df=91$，$p<0.001$）；信度检验得到总量表克隆巴赫Alpha系数为0.735。KMO值处于良好水平，球形检验拒绝了变量相互独立的原假设，说明题项间相关性显著、内部一致性良好，满足因子分析的前提条件，故此处不再重复报告检验细节。

\subsection{噪声影响的四维因子解构}

本研究以\textbf{主成分分析法}提取公因子，并采用凯撒正态化最大方差法进行旋转。旋转的目的是重新分配各因子所解释的方差，使因子载荷向0或1两极分化，从而让每个题项尽可能只在一个因子上呈现高载荷，便于因子命名与解释。

\subsubsection{公因子方差}

\begin{table}[H]
\centering
\caption{公因子方差}
\label{tab:communality}
\rowcolors{2}{tablight}{white}
\begin{tabular}{M{6cm}M{2.5cm}M{2.5cm}M{3cm}}
\bluetoprule
\rowcolor{tabhead}
\thead{题项} & \thead{初始值} & \thead{提取值} & \thead{提取程度} \\
\hline
交通噪声影响 & 1 & 0.539 & 中 \\
周围人生活噪声影响 & 1 & 0.501 & 中 \\
装修/施工噪声影响 & 1 & 0.750 & 高 \\
商业经营噪声影响 & 1 & 0.648 & 较高 \\
公共活动噪声影响 & 1 & 0.611 & 较高 \\
动物噪声影响 & 1 & 0.634 & 较高 \\
噪声入睡情况影响 & 1 & 0.616 & 较高 \\
噪声中断睡眠影响 & 1 & 0.572 & 中 \\
噪声注意力影响 & 1 & 0.742 & 高 \\
噪声身边关系影响 & 1 & 0.655 & 较高 \\
噪声精神状态影响 & 1 & 0.757 & 高 \\
噪声对身体健康影响 & 1 & 0.690 & 较高 \\
噪声对听觉影响 & 1 & 0.580 & 中 \\
噪声对极端情绪影响 & 1 & 0.668 & 较高 \\
\bluebottomrule
\end{tabular}
\end{table}

表\ref{tab:communality}给出各题项的公因子方差。初始值均为1，代表每个题项所包含的全部信息量；提取值表示公因子对该题项信息的解释比例，数值越大说明该题项的信息被公因子保留得越充分。结果显示，全部题项的提取值均大于0.5，处于中等及以上水平，其中“装修/施工噪声影响”“噪声精神状态影响”“噪声注意力影响”的提取值分别为0.750、0.757、0.742，提取程度为“高”；其余题项提取值介于0.501$\sim$0.690之间，提取程度为“中”或“较高”。所有题项的信息都能被公因子有效解释，无信息损失严重的题项，无需删除处理。

\begin{table}[H]
\centering
\caption{总方差解释}
\label{tab:variance}
\begin{subtable}[t]{0.49\textwidth}
\centering
\caption{提取载荷平方和}
\small
\setlength{\tabcolsep}{3pt}
\rowcolors{2}{tablight}{white}
\begin{tabular}{M{1.2cm}M{2.1cm}M{1.7cm}M{1.7cm}}
\bluetoprule
\rowcolor{tabhead}
\thead{成分} & \thead{可解释信息量} & \thead{方差百分比} & \thead{向上累积\%} \\
\hline
因子1 & 4.230 & 30.22\% & 30.22\% \\
因子2 & 1.950 & 13.93\% & 44.14\% \\
因子3 & 1.678 & 11.99\% & 56.13\% \\
因子4 & 1.105 & 7.89\% & 64.03\% \\
\bluebottomrule
\end{tabular}
\end{subtable}
\hfill
\begin{subtable}[t]{0.49\textwidth}
\centering
\caption{旋转载荷平方和}
\small
\setlength{\tabcolsep}{3pt}
\rowcolors{2}{tablight}{white}
\begin{tabular}{M{1.2cm}M{2.1cm}M{1.7cm}M{1.7cm}}
\bluetoprule
\rowcolor{tabhead}
\thead{成分} & \thead{可解释信息量} & \thead{方差百分比} & \thead{向上累积\%} \\
\hline
因子1 & 2.735 & 19.53\% & 19.53\% \\
因子2 & 2.622 & 18.73\% & 38.26\% \\
因子3 & 1.824 & 13.03\% & 51.29\% \\
因子4 & 1.783 & 12.73\% & 64.03\% \\
\bluebottomrule
\end{tabular}
\end{subtable}
\end{table}

表\ref{tab:variance}报告了公因子对原始信息的解释能力。旋转前，第一个因子单独解释了30.22\%的方差，旋转后其贡献率回落至19.53\%，而四个因子的方差贡献率分别为19.53\%、18.73\%、13.03\%和12.73\%，分布趋于均衡，说明旋转有效削弱了单一因子的主导地位。四个公因子的累积方差解释率为\textbf{64.03\%}，即原有14个题项六成以上的信息可由这四个因子承载，降维效果良好。

\subsubsection{公因子提取与命名}

\fig{图片17.png}{因子分析碎石图\label{fig:scree}}

公因子的提取数量结合碎石图拐点与特征值大小共同判定。碎石图以因子序号为横轴、对应特征值为纵轴，其中特征值代表单一公因子能够承载的原始变量信息量，其衰减规律可直观反映各因子的解释力度。

如图\ref{fig:scree}所示，折线在前4个因子区间内下降陡峭，自第5个因子起斜率大幅收窄、走势趋于平缓，曲线贴近横轴缓慢延伸，在第4、5因子之间形成明显的“肘部拐点”。数值上，前4个因子的特征值分别为4.230、1.950、1.678和1.105，均大于1；第5个因子特征值降至0.777，已低于1的临界阈值。碎石图拐点与特征值大于1准则的判定结果一致，不存在判定冲突，故提取4个公因子作为噪声影响的核心潜在维度；若继续纳入后续因子，其承载的信息量有限，多为随机测量误差，反而会降低模型的简洁性与可解释性。

\fig[1]{图片18.png}{噪声影响因子载荷热力图\label{fig:loading}}

结合旋转后的因子载荷热力图（图\ref{fig:loading}），各题项在其所属因子上的载荷绝对值均大于0.5，交叉载荷均小于0.3，因子结构清晰、区分度良好。为使各因子方向一致，本文按主载荷为正的方向调整因子得分符号，即因子得分越高表示该维度受到的噪声影响越强。依据各因子所涵盖题项的内容，将四个因子分别命名如下：

\textbf{（一）身心健康与认知干扰因子}

该因子包含噪声入睡情况影响、噪声注意力影响、噪声精神状态影响、噪声对听觉影响4个题项，主载荷介于0.720$\sim$0.865之间。其中入睡困难与听觉损害反映噪声造成的生理机能损伤，注意力涣散与精神状态波动反映噪声对认知与心理状态的干扰，两方面共同构成本因子所指的“身心健康与认知”损害。

\textbf{（二）情绪与社会关系影响因子}

该因子包含噪声中断睡眠影响、噪声身边关系影响、噪声对身体健康影响、噪声对极端情绪影响4个题项，主载荷介于0.736$\sim$0.804之间。该因子刻画的是噪声引发的情绪反应及其连带后果：极端情绪是核心表现，睡眠中断与身体不适是情绪应激在生理层面的反映，身边关系紧张则是负面情绪向人际互动外溢的结果。需要说明的是，“入睡困难”与“睡眠中断”分属两个因子，前者体现噪声对睡眠启动的直接干扰，后者则更多伴随情绪烦躁而出现，二者在成因上并不相同。

\textbf{（三）典型环境噪声源干扰因子}

该因子包含交通噪声影响、装修/施工噪声影响、公共活动噪声影响3个题项，主载荷介于0.677$\sim$0.843之间。该因子反映交通、施工、公共活动等公共性噪声源对居民的直接干扰，体现城市环境中典型噪声源的暴露水平与感知强度。

\textbf{（四）生活场景噪声源干扰因子}

该因子包含周围人生活噪声影响、商业经营噪声影响、动物噪声影响3个题项，主载荷介于0.633$\sim$0.791之间。该因子反映居住与生活场景中人际活动、商业经营、动物等贴近个体生活的噪声源带来的干扰，与第三因子所指的公共性噪声源形成区分。

\subsection{公因子得分计算}

为将题项层面的信息整合为可量化的维度指标，本研究依据因子得分系数矩阵计算每个样本在四个因子上的得分：

\begin{equation}
F_i=\sum_{j=1}^{m}\beta_{ij}\,ZX_j,\qquad i=1,2,3,4
\label{eq:factor_score}
\end{equation}

其中：$F_i$ 为第 $i$ 个公因子的得分；$m=14$ 为参与因子分析的题项总数；$\beta_{ij}$ 为因子得分系数矩阵中第 $i$ 个因子在第 $j$ 个题项上的系数；$ZX_j$ 为第 $j$ 个题项的标准化得分，其计算方式为

\begin{equation}
ZX_j=\frac{X_j-\overline{X}_j}{S_j},\qquad j=1,2,\dots,m
\label{eq:standardize}
\end{equation}

其中 $X_j$ 为第 $j$ 个题项的原始得分，$\overline{X}_j$ 与 $S_j$ 分别为该题项的样本均值与标准差。

式\eqref{eq:factor_score}中的求和范围覆盖全部14个题项，而非仅该因子载荷较高的题项，原因是因子得分系数矩阵同时利用题项间的相关结构对高分载荷题项的信息进行加权，对非主载荷题项也赋予较小的修正权重。由此得到的 $F_1\sim F_4$ 相互独立且量纲统一，可直接作为后续聚类分析的输入特征与回归模型中的解释变量。
